\documentclass[../main.tex]{subfiles}
\begin{document}
%==========================================TIÊU ĐỀ TẬP========================================
\begin{table}[h]
  \begin{adjustwidth}{-1cm}{}
    \begin{tabular}{>{\centering\arraybackslash}p{6cm}|>{\raggedright\arraybackslash}p{12.5cm}}
      \multirow{5}{6cm}
      % Cột bên trái: ÉPISODE và số tập
      {\\[-40pt]
        \flaregothic\fontsize{20pt}{20pt}\selectfont \centering
        {\contour{errorfive}{\textcolor{black}{ÉPISODE}}}\color{black}\\
        \burbank\fontsize{72pt}{72pt}\selectfont
        \includegraphics[height=3.12359cm]{../../../../Image/Book?/number/num36_alt.png}
        \\[18pt]
        % \flaregothic\fontsize{14pt}{14pt}\selectfont
        % \color{epattr}BIENVENUE, \uline{\color{Banpire} Mel} !
        % \flaregothic\fontsize{18pt}{18pt}\selectfont
        % \color{epattr}ÉPISODE PERDU
        \color{black}
      }
       &
      % Dòng thứ nhất: tên tập tiếng Nhật
      {\fotskip\fontsize{18pt}{18pt}\selectfont \ruby{知}{ち}\ruby{識}{しき}\ruby{合}{あ}\ruby{戦}{せん}}            \\[6pt]
      % Dòng thứ hai: tên tập tiếng Anh
       & {\cafeteria\fontsize{18pt}{18pt}\selectfont The War of Knowledge}                                 \\[4pt]
      % Dòng thứ ba: tên tập tiếng Việt
       & {\vnmsans\fontsize{18pt}{18pt}\selectfont Cuộc chiến tri thức}                                    \\[8pt]
      % Dòng thứ tư: nhân vật xuất hiện
       & {\caviar{\fontsize{14pt}{14pt}\selectfont Nhân vật: Tất cả các nhân vật xuất hiện từ đầu series}} \\[8pt]
      % \\[6pt]
      % Dòng cuối cùng: Thời gian diễn ra của tập (thường sẽ là ngày phát hành)
      % & {\continuum\fontsize{16pt}{16pt}\selectfont Ngày phát hành: 11/06/2022}\\[8pt]
    \end{tabular}
  \end{adjustwidth}
\end{table}
%===========================================NỘI DUNG==========================================
\color{black}

\begin{center}
  \uline{Vài ngày sau.}
  \pov{Hikawa Kuon}
\end{center}

Thời gian trôi qua nhanh hơn tôi tưởng. Mới đó mà đã gần một tuần kể từ khi Naomi chính thức được tái sinh ở thế giới này. Mọi chuyện diễn ra suôn sẻ một cách kỳ lạ. Dù vẫn còn chút ít sự tò mò và những lời xì xào, nhưng nhìn chung, sự hiện diện của cô bé mắt xanh đã trở thành một phần tất yếu của cuộc sống học đường thường nhật của lớp 1-C chúng tôi.

Một buổi sáng tầm giữa tháng Sáu, khi tôi đang mải mê kiểm tra lại đống tài liệu cho tiết tự học, cánh cửa lớp 1-C bỗng dưng bị đẩy mạnh ra.

``\textbf{Hỡi những kẻ phàm trần! Hãy lắng nghe thông điệp từ Vương quốc Bóng tối đây!}''

Shion lao vào lớp với vẻ mặt hớn hở, tay ôm khư khư một xấp tài liệu cũ kỹ có vẻ như được mượn từ thư viện cổ của thị trấn. Aku lếch thếch chạy theo sau, tay xách nách mang đủ thứ đồ dùng của câu lạc bộ Huyền bí, suýt chút nữa thì vấp ngã ngay ngưỡng cửa.

``Shion-san, có chuyện gì mà cậu hớt ha hớt hải vậy?'' tôi thở dài, gấp cuốn sách lại.

Cô nàng Phù thủy tự xưng không hề bận tâm tới câu hỏi, lao thẳng tới bàn của tôi và Naomi. Cô trải xấp tài liệu ra, đôi mắt tím sáng rực lên: ``Kuon! Nhìn này! Ta đã tìm thấy ghi chép về `Nghi thức Thanh tẩy Đêm trăng rằm'. Lễ hội sắp tới chính là thời điểm giao thoa mạnh mẽ nhất. Nếu chúng ta thực hiện nó, sự tồn tại của Naomi-chan sẽ hoàn toàn ổn định ở thế giới này!''

Cái tên ``Lễ hội Aogami'' vừa thốt ra khiến tôi khựng lại một nhịp. Ký ức về những ngày đầu tiên trôi dạt tới thế giới này, trong cái vòng lặp u ám của thị trấn, bỗng dưng ùa về rõ rệt. Tôi nhớ rất rõ, cũng vào một ngày nắng nóng như thế này, Hanabi và Suzune từng hào hứng rủ anh em tôi đi lễ hội. Lúc đó, cô nàng còn nửa đùa nửa thật dọa rằng nếu không tôn trọng Thần linh thì sẽ bị ``nghiệp quật'', còn Suzune thì sợ hãi cầu khấn không ngừng. Nhưng ở thời điểm đó, cái gọi là lễ hội mang một màu sắc bí ẩn và có chút gì đó đe dọa, khác hẳn với sự náo nhiệt chân thực của hiện tại.

``Nhắc mới nhớ, lễ hội năm nay chắc chắn sẽ khác hẳn hồi năm ngoái nhỉ?'' Nanase đứng cạnh lên tiếng, thu hút sự chú ý của cả nhóm. ``Tôi nghe nói các gian hàng đã bắt đầu được dựng lên ở khu vực đền rồi.''

``Đúng thế!'' Honoka chen vào, giọng đầy phấn khích. ``Nhớ năm ngoái thị trấn còn có mấy trò rút thăm trúng thưởng nữa mà. Không biết năm nay có không nhỉ?''

Naomi nghiêng đầu nhìn đống giấy tờ ngoằn ngoèo, rồi ngước nhìn Shion với vẻ mặt thản nhiên: ``Em cảm thấy mình đã rất ổn định rồi mà.''

``Không, không! Ngươi không hiểu được sự sâu xa của ma thuật đâu nhóc tì!'' Shion vung tay đầy kịch tính, rồi chỉ vào những dòng chữ Hán cổ trong tài liệu. ``Nhìn đây! Nghi thức Thanh tẩy Đêm trăng rằm là một đại điển cổ xưa, diễn ra vào đêm trăng tròn trùng với ngày hội. Mục đích của nó là dùng sự cộng hưởng linh khí của toàn bộ người dân thị trấn để xoa dịu các linh hồn vất vưởng. Với một người vừa `tái sinh' từ oán khí như ngươi, đây là cách tốt nhất để neo chặt linh hồn vào thực tại này, ngăn chặn việc bị `đào thải' bởi quy luật tự nhiên.''

Cô nàng nheo mắt, vẻ mặt trở nên cực kỳ nghiêm túc (dù vẫn mang đậm chất chuunibyou): ``Đây là sứ mệnh của Câu lạc bộ Huyền bí! Ta nhất định phải giúp ngươi trở thành một cư dân thực thụ của thế giới ánh sáng này!''

Không khí lớp 1-C lúc này rộn ràng hơn bao giờ hết. Những tiếng bàn tán về các gian hàng đồ ăn, các trò chơi dân gian và cả những bộ Yukata sặc sỡ lan tỏa khắp căn phòng. Ai nấy đều rạng rỡ, ánh mắt lấp lánh niềm vui khi đếm ngược từng ngày tới ngày hội lớn nhất của thị trấn Aogami. Thậm chí cả những học sinh vốn trầm tính nhất cũng bị cuốn vào vòng xoáy hào hứng này, tạo nên một sự cộng hưởng năng lượng tích cực mà tôi chưa từng thấy trước đây.

Tôi khẽ chạm vào mặt đồng hồ, cảm nhận một luồng điện nhẹ lướt qua da thịt. Trong khi lớp học đang xôn xao về những gian hàng, một giọng nói trầm thấp vang lên trong tâm trí tôi.

``\eng{Kuon-user, cậu nghe những gì mà cô ấy nói chứ?}'' Game lên tiếng, giọng ông ta mang một sắc thái nghiêm nghị hiếm thấy. ``\eng{Cái gọi là `Nghi thức Thanh tẩy Đêm trăng rằm' đó\dots\ nếu nó thực sự có thật từ thời cổ xưa như cô ta nói, nó có thể là tần số cộng hưởng mà chúng ta tìm kiếm. Nó là chìa khóa để mở cổng không gian.}''

Tôi khẽ nhíu mày, đáp lại bằng suy nghĩ: ``\eng{Ý ông là\dots\ nó có thể đưa chúng ta về nhà?}''

``\eng{Đúng vậy. Nhưng cũng có một vấn đề,}'' ông ta thở dài. ``\eng{Còn cô bé kia thì sao? Naomi là một linh hồn tái sinh từ tinh túy của thị trấn này. Nếu chúng ta rời đi, điểm neo giữ cô bé với thực tại này sẽ biến mất. Không có chúng ta, cô bé sẽ không còn ai để bám víu vào, và bóng tối có thể sẽ lại nuốt chửng cô bé.}''

Nỗi trăn trở ấy bóp nghẹt lấy trái tim tôi. Chúng tôi đã tìm thấy chìa khóa, nhưng cái giá phải trả có lẽ là sự cô độc của linh hồn nhỏ bé mà chúng tôi hứa sẽ bảo vệ. Chúng tôi phải rời đi, nhưng Naomi\dots\ em sẽ bám víu vào ai khi những người thân duy nhất của em biến mất khỏi thực tại này? Một phương án giải quyết, một phép màu khác\dots\ đó là thứ mà tôi và Game vẫn đang bế tắc trong việc tìm kiếm.

``\eng{\dots Đó thực sự là một vấn đề khá nan giải đấy. Tôi cũng đang nghĩ đến kịch bản đó đây\dots}'' tôi thở dài đáp lại.

``\eng{Dẫu cho lựa chọn là gì, nó có thể sẽ ảnh hưởng tới cô bé đấy,}'' Game nói thêm. ``\eng{Cậu và những người đi theo cậu sẽ cần thảo luận với nhau thật kỹ đó. Đặc biệt là với con bé Naomi kia.}''

``\eng{Tôi biết rồi.}''

Tiếng cười nói rộn ràng của cả lớp bỗng chốc bị cắt ngang bởi tiếng gõ thước khô khốc lên bảng đen. Yuko-sensei đã đứng ở bục giảng từ lúc nào, gương mặt cô hiện lên một nụ cười mà theo tôi thấy là cực kỳ đáng ngại.

``Tôi rất vui vì thấy các em có tinh thần hào hứng với Lễ hội như vậy,'' Yuko-sensei khoanh tay, đôi mắt sắc sảo nhìn cả lớp. ``Nhưng trước khi chạm được tới `vinh quang' của Lễ hội, có một rào cản nho nhỏ mà tất cả các em phải vượt qua.''

Một luồng khí lạnh lẽo bao trùm cả lớp. Chúng tôi đều biết điều gì sắp đến.

``Tuần sau chúng ta sẽ bước vào kỳ thi cuối kỳ,'' cô giáo dõng dạc thông báo. ``Và theo quy định của nhà trường, bất kỳ học sinh nào có điểm trung bình dưới mức yêu cầu sẽ phải ở lại trường học phụ đạo tăng cường suốt cả tuần Lễ hội.''

Thông tin ấy vừa được buông ra, cả lớp bỗng chốc lặng đi trong một tích tắc, rồi lập tức bùng nổ bởi những tiếng xì xào đầy hoảng hốt. Những gương mặt mới ban nãy còn rạng rỡ vì lễ hội, giờ đây đã tái mét như vừa nghe thấy bản án tử hình.

Tôi khẽ thở dài, tự rà soát lại mớ kiến thức trong đầu. Với tôi, đại số vốn là sở trường, những con số và phương trình luôn tuân theo một logic tuyệt đối mà tôi có thể dễ dàng làm chủ. Các môn học khác cũng ở mức ổn định, đủ để tôi không phải quá lo lắng về việc thi cử.

Thế nhưng, nhìn sang hai người đồng hành bên cạnh, tôi lại có chút quan ngại. Kaigou vốn là người của những hình khối và không gian, cô nàng có thể giải quyết các bài toán hình học hóc búa nhất - tôi biết điều đó từ bài kiểm tra mà Kuzuha đã từng cho chúng tôi làm hồi ở thế giới Nijisanji\footnote{Xem lại {\flaregothic JOUR \hbrk 8}, quyển Cầu Vồng Một.}. Suzuki thì lại là một ``phù thủy hóa học'' thực thụ, nhưng khi đứng trước các môn xã hội, em ấy trông chẳng khác gì một linh hồn đang lạc lối.

Tôi quay sang, hạ thấp giọng hỏi: ``\vie{\hl{Hai người tình hình thế nào? Có môn nào `báo động đỏ' không?}}''

Kaigou nhún vai, vẻ mặt thản nhiên như thường lệ: ``\vie{\hl{Hình học thì ổn, còn lại chắc là\dots{} đủ để qua môn.}}''

Suzuki cũng cười gượng, gãi đầu: ``\vie{\hl{Hóa học thì em cân được, nhưng mấy môn khác thì chắc cũng chỉ hy vọng chạm mức trung bình thôi ạ.}}''

Trong khi chúng tôi còn đang lo cho bản thân, thì ở phía bàn của Suzune, thảm họa thực sự đã bắt đầu.

``{\Large\textbf{KHÔNGGGGGG!!!}}''

Cô nàng linh vật của lớp gục mặt xuống bàn, mái tóc vàng-cam rũ xuống che hết cả mặt. ``Học phụ đạo vào đúng tuần có gian hàng bánh bao ngọt sao!? Đây là tận thế! Chính là tận thế rồi!''

Yae và Honoka cũng không khá khẩm hơn. Yae đập đầu xuống bàn than vãn về việc không được tham gia giải đấu thể thao truyền thống, còn Honoka thì lẩm bẩm điều gì đó về việc bỏ lỡ cơ hội buôn chuyện với toàn thị trấn.

Yuko-sensei mỉm cười hài lòng trước sự suy sụp của hội ``báo động đỏ'', rồi cô quay sang phía tôi, Sakura, Touka và Nanase - những người có thành tích học tập tốt nhất lớp.

``Kuon-san, Shinomiya-san, Yukihara-san, Furukawa-san. Tôi hy vọng các em có thể hỗ trợ các bạn một chút. Chúng ta không muốn lớp 1-C vắng mặt bất kỳ ai trong ngày hội lớn nhất năm, đúng không?''

Trong khi tôi đang định lên tiếng đồng ý, Touka đã khoanh tay, đôi mắt vàng của cô nheo lại nhìn Suzune vẫn đang nằm bẹp trên bàn: ``Suzune-san, mình đã bảo cậu từ đầu năm học rồi mà. Nếu cứ mải mê với bánh bao mà bỏ bê bài vở, sớm muộn gì ngày này cũng tới thôi.''

``Ư\dots\ Touka-chan ác quá\dots'' cô nàng ``Linh vật'' thút thít, không dám ngước mặt lên.

Bất chợt, Shion đang đứng thẳng lưng, tay chống nạnh cười đầy tự phụ: ``Hừm! Phù thủy Bóng tối như ta không cần lo lắng về mấy con chữ tầm thường đó. Kiến thức của ta vượt xa sự hiểu biết của nhân loại!''

Aku đứng bên cạnh khẽ kéo áo cô bạn, giọng nhỏ nhẹ nhưng đầy tính sát thương: ``Nhưng Shion-chan\dots\ bài kiểm tra Sinh học hôm nọ, bà chỉ được 15 điểm thôi. Tại bà mải vẽ vòng tròn ma thuật vào mặt sau tờ giấy thi đó.''

Nụ cười trên môi cô nàng Phù thủy đông cứng lại. Tôi ngay lập tức nắm bắt cơ hội, mỉm cười hỏi: ``Vậy sao? Shion-san, cậu trả lời cho tôi xem cấu tạo của ty thể trong tế bào gồm những thành phần nào?''

``Cái đó\dots\ là một bí thuật không thể tiết lộ!'' Shion khựng lại, đôi mắt tím đảo liên tục, giọng nói bắt đầu ấp úng. ``Nó\dots\ nó gồm có lớp vỏ bóng tối và\dots\ nhân của sự hủy diệt?''

Tôi thở dài. ``Được rồi, câu hỏi dễ hơn nhé. Điều kiện xác định của hàm số $x^2 + 8$ là gì?''

``Điều kiện\dots\ xác định? Ý của ngươi là xác định xem ta nên dùng ma thuật nào sao?'' Shion nghiêng đầu thắc mắc.

\dots Vậy là đã rõ. Danh sách ``báo động đỏ'' chính thức có thêm cái tên của cô nàng Phù thủy tự xưng này.

``Này, sao mấy người học hành kiểu gì vậy?'' Yae ngẩng đầu lên, vẻ mặt đầy thắc mắc. ``Nhìn Yuka-san kìa, cậu ấy ngủ suốt cả ngày trong lớp mà điểm lúc nào cũng cao. Còn Yuki-san nữa, cô ấy suốt ngày đi dạo bên ngoài, có bao giờ thấy cầm sách đâu mà kết quả vẫn ổn định?''

Nàng bờm Mèo nghe thấy tên mình, khẽ hé một mắt ra khỏi giấc ngủ ngái ngủ, miệng cười tinh quái: ``Hehe, đó là vì mình có năng lượng ngoại cảm mừ\~{} Chỉ cần nằm đó, tri thức của giáo viên sẽ tự động truyền vào đầu mình thôi\~{}''

Thực ra, tôi biết lý do không phải vậy. Yuka có khả năng tiếp thu và xử lý thông tin cực nhanh; cô ấy chỉ cần nghe lướt qua một lần là đã nắm vững bản chất vấn đề. Còn Yuki, cô nàng này có một óc quan sát và trí nhớ hình ảnh đáng kinh ngạc. Những gì cô ấy thấy trên đường phố, trong thiên nhiên đều được cô ấy liên hệ trực tiếp với kiến thức địa lý và xã hội học một cách tự nhiên.

``Đó gọi là thiên phú rồi, bọn mình không so được với họ đâu, Yae-san,'' Suzuki cười khổ đáp lời.

Tôi nhìn quanh lớp, thấy rõ sự khác biệt giữa những người học dễ và học khó. Ánh mắt tôi dừng lại ở Naomi đang lặng lẽ quan sát, rồi quay lại nhóm ``báo động đỏ'' đang rầu rĩ. Một ý tưởng chợt lóe lên trong đầu: chúng tôi không chỉ cần học, mà cần học cùng nhau. Đó có lẽ là cách khả dĩ nhất để vượt qua khó khăn này.

``Vâng, thưa cô. Chúng em sẽ lo liệu việc này,'' tôi dõng dạc đáp lời. Touka chỉ khẽ gật đầu, còn Nanase thì giơ nắm đấm đầy tự tin. Sakura mỉm cười dịu dàng, ánh mắt lấp lánh sự nhiệt huyết: ``Cô cứ yên tâm ạ, chúng em sẽ không để ai phải ở lại đâu!''

Yuko-sensei nhìn chúng tôi với nụ cười đầy ẩn ý: ``Tốt lắm. Đừng để cô phải thất vọng nhé, đặc biệt là với môn Sinh học mà cô trực tiếp giảng dạy đấy. Cô mong chờ kết quả của các em.''

Vậy là nhiệm vụ bất đắc dĩ của chúng tôi bắt đầu. Cả lớp 1-C biến thành một ``lò luyện thi'' cấp tốc ngay tại lớp sau giờ học. Tôi đứng dậy, gõ nhẹ lên bàn để thu hút sự chú ý của nhóm chủ chốt.

``Để đảm bảo không ai bị bỏ lại, chúng ta cần phân công gia sư,'' tôi dõng dạc nói. ``Touka-san, cậu là người đứng đầu khối, phiền cậu đảm nhận phần lớn các môn học nhé?''

``Sẵn sàng thôi,'' Touka đẩy gọng kính vô hình, ánh mắt kiên định. ``Mình sẽ đảm bảo các cậu ấy nắm vững kiến thức trọng tâm.''

``Nanase-san, cậu kèm cặp môn Văn học và Lịch sử nhé? Cách cậu kể chuyện rất lôi cuốn, chắc chắn sẽ giúp mọi người dễ nhớ hơn,'' tôi quay sang cô nàng tóc xanh.

``Cứ tin ở tôi! Tôi sẽ biến mấy trang sử khô khan thành kịch bản phim hành động luôn!'' Nanase giơ ngón tay cái đầy tự tin.

``Tôi sẽ phụ trách môn Sinh học,'' Sakura lên tiếng, giọng nói tràn đầy năng lượng. ``Tôi đã chuẩn bị sẵn một xấp tài liệu về các hệ sinh thái và di truyền học rồi.''

Tôi nhìn sang hai người bạn thân của mình: ``Kaigou, cậu phụ trách phần Hình học, còn Suzuki lo môn Hóa học cho cả lớp nhé. Tôi sẽ chạy đôn chạy đáo hỗ trợ môn Đại số cho nhóm `báo động đỏ'. Mọi người đồng ý chứ?''

Tất cả đều gật đầu tán thành. Một ``chiến dịch cứu trợ'' quy mô lớn chính thức được kích hoạt.

Cô bé mắt xanh ngồi im lặng bên cạnh Suzuki, đôi mắt lục bảo dõi theo từng con số, từng dòng chữ trên những trang sách. Naomi không cầm bút, cũng không làm bài, em chỉ đơn giản là lắng nghe. ``\vie{Mọi người đang chiến đấu với những con chữ sao?}'' Naomi khẽ hỏi khi thấy Yae đang vò đầu bứt tai bên cạnh Nanase.

``\vie{Đúng vậy đấy Naomi-chan,}'' Suzuki cười khổ, lau mồ hôi trên trán. ``\vie{Nếu không thắng được chúng, các anh chị sẽ không được đi chơi Lễ hội đâu.}''

Naomi gật đầu vẻ thấu hiểu. Em vươn bàn tay nhỏ nhắn ra, khẽ chạm vào trang sách của Suzuki. Một luồng khí mát lành, dịu nhẹ tỏa ra từ đầu ngón tay em, lan tỏa khắp căn phòng, xua tan đi cái nóng bức và sự căng thẳng đang đè nặng lên vai mọi người.

``\vie{Em sẽ hỗ trợ mọi người theo cách của em,}'' Naomi thì thầm, đôi mắt khẽ nhắm lại. ``\vie{Em sẽ giữ cho tâm trí mọi người luôn tĩnh lặng như mặt hồ, để những tri thức này có thể đọng lại một cách dễ dàng nhất ạ.}''

Và thế là, chiến dịch ``Vượt qua tường thành thi cử chạm tới ngày lễ hội nhộn nhịp'' của lớp 1-C chính thức bắt đầu.

\ct

Phòng học lớp 1-C vào những buổi chiều muộn giờ đây chẳng khác gì một chiến trường thực sự, với vũ khí của chúng tôi là bút mực, thước kẻ và những chồng sách giáo khoa cao ngất ngưởng.

Ở góc cửa sổ, Sakura đang đứng giữa một vòng tròn các bạn học, trong đó có cả Hotaru và Mitsuki. Cô nàng mỉm cười rạng rỡ, tay cầm một mô hình cấu trúc DNA tự chế bằng những hạt cườm đầy màu sắc.

``Mọi người nhìn này! Nếu chúng ta coi các liên kết hydro giống như những cái nắm tay của các cặp đôi đi hội quân, thì A luôn đi với T và G luôn đi với X\footnote{A, T, G, X lần lượt là adenine, thymine, guamine và cytosine. Chúng là bốn loại axit nucleic, được dùng để mã hóa các thông tin di truyền, được lưu trữ trong DNA. Trong các tài liệu nước ngoài, chữ C thường được sử dụng cho cystosine thay vì chữ X ở Việt Nam.}. Nếu nắm nhầm tay là sẽ xảy ra `scandal' di truyền ngay lập tức!''

Cách ví von sinh động và đầy hình tượng của cô nàng Vu nữ khiến môn Sinh học vốn khô khan trở nên gần gũi vô cùng. Cô nàng không ép mọi người học thuộc lòng định nghĩa, mà dẫn dắt họ bằng những câu kể về sự kỳ diệu của sự sống. Mitsuki, người vốn hay lơ đãng, giờ đây cũng đang chăm chú ghi chép, đôi mắt mở to đầy hứng khởi.

Trong khi đó, ở khu vực bảng đen, Nanase đang múa may quay cuồng với những viên phấn màu. Cô ấy không dạy Lịch sử bằng những con số khô khan, mà biến mỗi triều đại, mỗi cuộc chiến thành một vở kịch drama kịch tính. Không biết có phải do cô xem phim nhiều không, nhưng nó vẫn hiệu quả ở một khía cạnh nào đó, khi cách tiếp cận bài học kiểu này cũng mới lạ.

``Và ngay lúc này! Vị tướng quân này đã tung ra một đòn quyết định, thay đổi hoàn toàn cục diện chính trị!'' cô Tiểu thư vung tay đầy kịch tính, khiến nhóm của Honoka và Yae nghe như nuốt từng lời. ``Yae-san, nếu cậu là vị tướng đó, cậu sẽ làm gì để bảo vệ gian hàng thể thao của mình?''

``Tớ sẽ\dots{} tớ sẽ xông thẳng vào trung quân địch!'' Yae hét lên đầy nhiệt huyết, dường như đã hoàn toàn quên đi nỗi sợ hãi môn Lịch sử. Kỹ năng sư phạm của Nanase chính là biến học sinh thành nhân vật chính trong dòng chảy của thời gian, khiến họ cảm nhận được sức nặng của quá khứ thông qua những liên hệ thực tế.

Ngược lại hoàn toàn với không khí sôi nổi đó, ở dãy bàn giữa, Touka đang duy trì một sự tĩnh lặng đáng sợ. Cô ngồi đối diện với Suzune và Shion, đôi mắt sắc bén của nàng Thần đồng không bỏ sót một biểu cảm nào trên gương mặt họ.

``Suzune-san, đừng chỉ nhìn vào con số, hãy nhìn vào bản chất của hàm số này.'' Giọng Touka điềm tĩnh nhưng vô cùng quyết đoán. Cô nàng yêu cầu sự chuẩn xác tuyệt đối và phương pháp giải bài gãy gọn. ``Shion-san, nếu cậu còn tiếp tục nhầm lẫn giữa $\sin$ và $\cos$, mình sẽ bắt cậu chép phạt công thức bằng chính cái mực `bóng tối' của cậu đấy.''

Cô nàng Phù thủy run rẩy, lập tức cúi đầu làm bài không dám ho he. Phương pháp của cô nàng Thần đồng là ``thiết quân luật'' tri thức; cô ép đối phương phải đi sâu vào bản chất cốt lõi của vấn đề thay vì chỉ học vẹt công thức. Nó mang lại hiệu quả cực cao đối với những người có tư duy hay nhảy cóc như Shion. Nhưng với Suzune - cô nàng lúc nào cũng để óc trên mây\dots\ tôi không chắc liệu có tác dụng hay không.

Tôi đứng ở giữa lớp, đóng vai trò như một người điều phối tổng thể. Tôi chạy đôn chạy đáo từ nhóm này sang nhóm khác, giải đáp những thắc mắc vụn vặt và cung cấp các mẹo làm bài trắc nghiệm nhanh mà tôi đúc kết được từ mười hai năm đèn sách thời trung học và hơn hai năm rưỡi dùi mài kinh sử học đại học ở thế giới thực. Kể cả cái mẹo ``chọn toàn C'' của ông thầy nào đó trên mạng, dẫu tôi chưa bao giờ dám thử.

``Nhìn này, với dạng bài Đại số này, chúng ta có thể loại trừ phương án sai dựa vào điều kiện xác định trước,'' tôi vừa hướng dẫn cho Mari, vừa để mắt tới sự tiến triển của cả lớp. ``Mọi người làm tốt lắm! Cố gắng lên nào!''

Ngay cả các bạn học khác như Azuki hay Inari cũng không đứng ngoài cuộc. Họ tự giác lập thành những nhóm nhỏ, trao đổi bài vở và hỗ trợ lẫn nhau những phần kiến thức mà họ nắm vững.

Ở một góc bàn yên tĩnh phía sau, Naomi ngồi lặng lẽ bên cạnh Suzuki. Cô bé không tham gia vào những cuộc thảo luận sôi nổi, đôi mắt lục bảo trong veo chỉ lặng lẽ dõi theo những gương mặt đang miệt mài bên trang sách. Naomi quan sát cách Sakura cười, cách Touka nghiêm nghị, và cách Nanase hăng say kể chuyện. Đôi khi em nghiêng đầu như thể đang cố gắng thấu hiểu cái gọi là ``áp lực thi cử'' - một khái niệm hoàn toàn xa lạ với một linh hồn đã ngủ yên ngàn năm.

Suzuki khẽ mỉm cười nhìn cô bé, tay vẫn không ngừng chấm bài tập hóa học cho một cậu bạn ngồi cạnh: ``\vie{Em thấy thế nào, Naomi-chan? Không khí này\dots{} có chút lạ lẫm với em phải không?}''

Naomi khẽ gật đầu, môi mím lại thành một nụ cười nhạt: ``\vie{Mọi người\dots{} đều đang tỏa sáng một cách kỳ lạ.}''

\ct

Cách đó không xa, Kaigou đang đối mặt với một nhóm học sinh đang ``đau đầu'' vì môn Hình học không gian. Với kinh nghiệm từng kèm cặp Suzuki từ nhỏ, cô nàng có một phong thái sư phạm vô cùng điềm tĩnh và thực tế. Thay vì chỉ dùng phấn vẽ lên bảng, cô ấy tận dụng mọi đồ vật xung quanh - từ hộp bút, thước kẻ đến cả những chồng sách - để dựng lên các mô hình 3D ngay trên mặt bàn.

``Đừng cố tưởng tượng trong đầu nếu chưa thấy nó,'' Kaigou nhẹ nhàng xoay cái hộp bút để minh họa cho một mặt phẳng cắt qua khối lăng trụ. ``Hình học là sự sắp xếp của không gian. Chỉ cần hiểu được vị trí của các điểm, đường thẳng sẽ tự hiện ra.''

Sự kiên nhẫn và cách tiếp cận trực quan của Kaigou khiến những khái niệm trừu tượng nhất cũng trở nên dễ hiểu. Nhóm bạn học xung quanh gật đầu lia lịa, những ánh mắt bối rối dần được thay thế bằng sự minh mẫn.

Trong khi đó, ``phù thủy hóa học'' Suzuki lại biến góc lớp thành một buổi trình diễn khoa học thú vị. Cô nàng không mang theo hóa chất thật vào lớp, nhưng cách em mô tả các phản ứng lại sống động đến mức mọi người như ngửi thấy mùi đặc trưng của các chất khí.

``Hóa học không chỉ là những ký hiệu vô hồn trên trang giấy đâu!'' em ấy cười tươi, tay múa may minh họa cho sự trao đổi electron. ``Hãy coi axit và bazơ như hai người bạn đang trao đổi quà cho nhau để đạt đến trạng thái cân bằng. Chỉ cần hiểu được tâm tính của từng nguyên tố, việc cân bằng phương trình sẽ dễ như ăn bánh bao vậy!''

Sự nhiệt huyết tuổi trẻ và cách giải thích hóm hỉnh của Suzuki truyền lửa cho rất nhiều bạn học vốn sợ những con số lẻ và hóa trị phức tạp.

Tôi khẽ mỉm cười, cảm nhận được sự cộng hưởng mạnh mẽ của tri thức đang lan tỏa khắp phòng học. Naomi, sau một lúc ngồi cạnh Suzuki, giờ đây đã lặng lẽ di chuyển sang cạnh tôi. Cô bé ngồi xếp bằng trên ghế, đôi tay nhỏ nhắn đặt lên đầu gối, đôi mắt lục bảo dõi theo từng chuyển động của các ``gia sư''.

Em không tham gia vào những cuộc thảo luận sôi nổi, em cũng không làm bài, em chỉ đơn giản là hiện diện ở đó như một linh hồn tĩnh lặng giữa cơn bão. Đôi khi, cô bé lại ngước nhìn tôi, như muốn tìm kiếm một lời giải thích cho sự nỗ lực phi thường mà con người đang bỏ ra chỉ vì một kỳ thi.

``\vie{Onii-chan,}'' Naomi thì thầm, giọng nói trong trẻo như tiếng chuông gió. ``\vie{Mọi người đều đang cố gắng hết sức vì niềm vui của Lễ hội sao?}''

Tôi gật đầu, đặt tay lên mái tóc mềm của em: ``\vie{Đúng vậy, Naomi-chan. Đôi khi chúng ta phải vượt qua những bức tường cao nhất để có thể chạm tay vào niềm hạnh phúc thực sự.}''

Naomi khẽ nhắm mắt lại, một nụ cười nhạt hiện trên môi. Em không cầm bút, em cũng không làm bài, nhưng sự hiện diện bình yên của em cạnh tôi dường như đang tiếp thêm sức mạnh cho cả lớp 1-C.

Quả thực, trong căn phòng đầy áp lực này, ánh sáng của tri thức và sự quyết tâm đang bùng cháy rực rỡ hơn bao giờ hết, báo hiệu cho một cuộc lội ngược dòng ngoạn mục của lớp 1-C.

\ct

Nhưng rồi, những trắc trở trong việc ôn thi của chúng tôi đã bắt đầu lộ ra.

Sự nhiệt huyết ban đầu rồi cũng dần bị bào mòn bởi sức nặng của thời gian và khối lượng kiến thức khổng lồ. Khi kim đồng hồ nhích dần tới con số bảy giờ tối của ngày ôn tập thứ tư, bầu không khí trong lớp 1-C bắt đầu nhuốm màu mệt mỏi rã rời. Kết hợp với thời tiết bí bức của đầu hè càng khiến cho chúng tôi nản lòng hơn.

``Bánh bao\dots{} bánh bao nhân thịt\dots{} đừng chạy mà\dots{}''  Tiếng lẩm bẩm mê sảng của Suzune vang lên giữa sự tĩnh lặng. Cô nàng đã gục mặt xuống chồng bài tập đại số từ lúc nào, nước dãi chảy cả ra mép vở.

Bên cạnh đó, Shion cũng chẳng khá khẩm hơn. Mái tóc bạch kim thường ngày được chải chuốt kỹ lưỡng giờ đây rũ rượi, đôi mắt quầng thâm như thể vừa trải qua một nghi lễ hắc ám kéo dài mười ngày mười đêm.

``Sin bình phương cộng cos bình phương ($\sin^2\theta+\cos^2\theta$) bằng\dots{} bằng một cơn ác mộng bóng tối\dots{}'' nàng Phù thủy tự xưng lẩm bẩm, tay vẫn cầm cây bút vẽ linh tinh những vòng tròn vô nghĩa lên mặt bàn.

Yae thì nằm bò ra ghế, hai tay buông thõng, đôi mắt lờ đờ nhìn lên trần nhà. Với một người quen vận động như cậu ấy, việc phải ngồi yên một chỗ và nhồi nhét những con số là một cực hình thực sự. Ngay cả Honoka, người vốn có nguồn năng lượng ``tám chuyện'' vô tận, giờ đây cũng chỉ biết thở dài thườn thượt, tay chống cằm nhìn đống tài liệu lịch sử của Nanase như nhìn vào một mê cung không lối thoát.

Thấy tình hình có vẻ trầm trọng, Touka đành mềm lòng, dù gương mặt vẫn giữ vẻ nghiêm nghị nhưng giọng nói đã có phần dịu đi: ``Mọi người, đừng bỏ cuộc lúc này. Chúng ta chỉ còn ba ngày nữa là tới kỳ thi rồi. Uống chút nước và nghỉ ngơi năm phút đi.''

Sakura cũng đi vòng quanh, đặt những viên kẹo ngọt lên bàn của từng người kèm theo một nụ cười cổ vũ: ``Cố lên nào mọi người! Nghĩ đến cảnh chúng ta cùng nhau đi dạo dưới pháo hoa lễ hội đi, mọi mệt mỏi sẽ tan biến ngay thôi!''

Nanase thì cố gắng khuấy động bầu không khí bằng một câu chuyện lịch sử về sự kiên trì của các chiến binh cổ đại, nhưng có vẻ như ``khán giả'' của cô nàng đã quá kiệt quệ để có thể hào hứng như những ngày đầu.

Tôi đứng tựa lưng vào bục giảng, lòng trĩu nặng sự lo âu. Nhìn những gương mặt phờ phạc kia, tôi biết rằng giới hạn chịu đựng của họ đã chạm đáy. Nó làm tôi liên tưởng tới hồi tôi ôn thi vào cấp 3, cả lớp của tôi hồi đó không khác nào một bầy ``xác sống'' nhất là từ những kẻ không học bài từ đầu năm. Nếu cứ tiếp tục ép buộc như thế này, bộ não của họ sẽ ``cháy khét'' trước khi kịp bước vào phòng thi.

Shino từ phía bàn của mình bước lại gần chỗ chúng tôi, gương mặt cô ấy cũng không giấu nổi vẻ mệt mỏi nhưng vẫn còn tỉnh táo hơn nhiều so với nhóm Suzune. Cô nhìn Suzune đang ``bơi'' trong đống giấy tờ rồi khẽ thở dài: ``Tình hình này có vẻ nghiêm trọng hơn mình tưởng. Nếu cứ tiếp tục thế này, ngay cả phép màu của Câu lạc bộ Huyền bí cũng chẳng cứu nổi điểm số của họ đâu.''

Kaigou khoanh tay, đôi mắt xanh lam của cô hiện lên sự ngán ngẩm: ``Chúng tôi cũng đang rất lo. Nhưng dường như giới hạn của họ đã chạm tới mức đỏ rồi. Ép thêm nữa sẽ chỉ phản tác dụng.''

``Mọi người đã nỗ lực rất nhiều rồi mà\dots'' Suzuki nhìn sang Naomi, bàn tay khẽ nắm chặt.

Shino nhìn theo ánh mắt của Suzuki, cô khẽ hạ thấp giọng đủ để chỉ bốn người chúng tôi nghe thấy: ``Naomi-chan có vẻ rất quan tâm đến các anh chị của mình nhỉ? Đôi mắt của cô bé\dots{} quả thật đang chứa đựng một điều gì đó rất thuần khiết.''

Tôi quay sang nhìn Naomi đang ngồi cạnh mình, em vẫn giữ vẻ điềm tĩnh đến lạ thường, đôi mắt lục bảo dõi theo sự rã rời của cả lớp với một sự trăn trở kín đáo. Đúng lúc ấy, cô bé mắt xanh bỗng đứng dậy khỏi ghế. Không một lời báo trước, em bước ra giữa lớp học đang chìm trong sự rã rời.

Em bước tới trước bục giảng, đôi tay nhỏ nhắn đan vào nhau trước ngực. Một vầng sáng xanh dịu nhẹ, mỏng manh như sương sớm bắt đầu tỏa ra từ cơ thể nhỏ bé của em, lan nhanh khắp phòng học như một làn gió mát lành thổi qua sa mạc.

Ngay lập tức, bầu không khí bí bức của buổi chiều đầu hè bỗng chốc tan biến. Tiếng thở dốc mệt mỏi lặng đi, thay vào đó là một sự tĩnh lặng tuyệt đối nhưng vô cùng dễ chịu. Suzune bỗng giật mình tỉnh giấc, đôi mắt xanh lá chớp chớp, không còn vẻ lờ đờ mà thay vào đó là một sự minh mẫn lạ thường.

``Ơ\dots{} đầu mình nhẹ bẫng thế này?'' Suzune thốt lên, tay đưa lên xoa thái dương. ``Mọi mệt mỏi biến đâu mất tiêu rồi!''

Shion cũng bật dậy, mái tóc bạch kim như được tiếp thêm sinh lực: ``Hừm! Ta cảm nhận được một nguồn linh khí tinh khiết đang chảy trong huyết quản! Phép thuật của ta\dots{} ý ta là trí não của ta đã sẵn sàng chinh phục mọi đỉnh cao!''

Bầu không khí học tập sau đó thay đổi hoàn toàn. Không còn những tiếng than vãn hay gục mặt xuống bàn, cả lớp 1-C hăng say lao vào đống tài liệu với một sự tập trung cao độ mà chính tôi cũng phải kinh ngạc. Những con số, công thức và sự kiện lịch sử giờ đây dường như tự động sắp xếp vào bộ não của họ một cách trật tự dưới sự hỗ trợ vô hình từ phép màu của Naomi.

``Mọi người thấy sao? Mình đã phổ nhạc toàn bộ các công thức Hóa học của Suzuki-san thành một bản ballad rồi này!'' Azuki hào hứng khoe, tay gõ nhịp lên mặt bàn.

Mari thì giơ lên những trang vở rực rỡ sắc màu: ``Còn mình thì dùng sơ đồ tư duy kết hợp với minh họa vẽ tay. Nhìn vào đây là thấy cả một hệ sinh thái của Sakura-san hiện ra luôn!''

Ayame khẽ gật đầu, đôi mắt dịu dàng nhìn về phía Naomi: ``Lạ thật đấy, mỗi khi nhìn thấy Naomi-chan, tôi lại cảm thấy tâm trí mình tĩnh lặng đến mức có thể đọc thông viết thạo cả những đoạn văn cổ khó nhất.''

Hotaru đứng cạnh cũng tán thành: ``Đúng thế, năng lượng của cô bé thật kỳ diệu. Mình cảm giác như mình có thể thức trắng đêm mà vẫn tỉnh táo vậy.''

Hanabi thì vẫn nhiệt huyết như mọi khi, cô nàng đập tay xuống bàn đầy quyết tâm: ``Được rồi! Với đà này, lớp 1-C chúng ta chắc chắn sẽ quét sạch mọi đề thi để tiến thẳng tới Lễ hội!''

\ct

Không chỉ dừng lại ở lớp học, một số bạn học như Inari, Yae và Honoka còn nỗ lực đến mức xin phép được sang nhà Hikawa chúng tôi để vừa học nhóm, vừa tranh thủ được ``cưng nựng'' và chơi đùa cùng Naomi.

Yae vừa cười vừa ôm chầm lấy cô bé, xoa xoa đôi má phúng phính của cô bé: ``Naomi-chan giỏi quá đi mất, em đúng là thiên thần cứu mạng của bọn chị đấy!'' Em chỉ mỉm cười hiền hậu, ngoan ngoãn để các chị ``vần vò'' mà không hề phàn nàn.

Inari, với tư cách lớp trưởng, khẽ điều chỉnh lại chiếc nơ trên cổ áo, ánh mắt nghiêm nghị nhưng chứa chan sự cảm kích: ``Mọi người, đừng quên rằng Naomi-chan đã vất vả thế nào để giúp chúng ta. Chúng ta nhất định phải đạt kết quả tốt để không phụ lòng cô bé.''

Honoka thì cười hì hì, tay làm dấu chiến thắng: ``Yên tâm đi Lớp trưởng! Tớ đã thoát khỏi kiếp `xác sống' rồi. Lần này tớ sẽ thi cho Yuko-sensei phải lác mắt nhìn luôn!''

Tuy nhiên, Suzuki không khỏi bày tỏ sự lo lắng khi thấy cô em gái bé nhỏ liên tục phải sử dụng năng lượng để thanh lọc tâm trí cho cả lớp. Trong một buổi tối muộn sau khi các bạn đã về hết, em khẽ hỏi Naomi khi đang giúp cô bé chải tóc:

``\vie{Naomi-chan, em làm thế này có mệt lắm không? Chị thấy em bị mọi người `nhờ vả' hơi nhiều rồi đấy. Em không cần phải cố gắng đến thế đâu.}''

Naomi ngước nhìn Suzuki qua tấm gương, đôi mắt lục bảo lấp lánh sự kiên định: ``\vie{Không sao đâu Suzuki-onee-chan. Với em, đây là cách tốt nhất để em có thể giúp sức cho các anh chị vượt qua `tường thành' này. Được nhìn thấy mọi người vui vẻ học tập và cùng nhau hướng tới Lễ hội, em cảm thấy mình cũng đang thực sự sống, thực sự là một phần của gia đình này ạ.}''

``\vie{Với em, được ở bên cạnh và ủng hộ mọi người là hạnh phúc lớn nhất,}'' cô bé tiếp lời, bàn tay nhỏ nhắn nắm lấy tay Suzuki. ``\vie{Dù sau này có chuyện gì xảy ra, dù chúng ta có đi tới đâu, em vẫn muốn được sát cánh cùng anh chị đến tận cùng ạ.}''

Suzuki lặng người, đôi mắt thoáng chút bối rối nhưng rồi lập tức được thay thế bằng một sự xúc động nghẹn ngào. Em ôm chặt lấy Naomi vào lòng, như muốn che chở cho linh hồn bé nhỏ ấy khỏi mọi giông bão của định mệnh.

Đứng ở ngoài cửa, tôi cảm thấy tim mình thắt lại. Câu nói ``nguyện theo tới cùng'' của Naomi, theo tôi, không chỉ đơn thuần là lời hứa của một đứa trẻ, mà dường như là một lời tiên tri cho cuộc hành trình xa xôi mà chúng tôi sắp phải đối mặt. Sự gắn kết này, tình cảm này\dots{} nó đã vượt xa khỏi ranh giới của một nhiệm vụ bảo vệ đơn thuần.

Lời nói chân thành của cô bé khiến Suzuki lặng người vì xúc động, còn tôi và Kaigou đứng ở cửa phòng chỉ biết nhìn nhau, mỉm cười đầy tự hào về linh hồn bé nhỏ mà chúng tôi đã mang về từ bóng tối.

\ct

\pov{Hikawa Suzuki}

Ngày thi vượt vũ môn cuối cùng cũng đến.

Cả ngôi trường Aogami như chìm vào một bầu không khí trang nghiêm đến lạ kỳ. Tiếng loa thông báo vang vọng khắp các hành lang, xen lẫn tiếng lật giấy sột soạt và tiếng bút chạy hối hả trên mặt đề thi, những âm thanh chủ đạo của một cuộc chiến không tiếng súng.

Bước vào phòng thi của lớp 1-C, tôi có thể cảm nhận rõ một luồng điện tích cực đang chạy qua từng dãy bàn. Không còn những tiếng thở dài hay những cái gục đầu bất lực như mấy ngày trước. Thay vào đó, tôi nhìn thấy những ánh mắt đang rực cháy quyết tâm của các bạn.

Suzune, người vốn thường xuyên thất thần, giờ đây đang nhìn chằm chằm vào đề bài Toán với một sự tập trung cao độ mà tôi chưa từng thấy. Ánh mắt cô ấy sắc sảo lạ thường, đôi môi mím chặt, tay cầm bút vạch ra những dòng giải toán một cách dứt khoát. Tôi thầm nhủ: ``\textit{Cố lên nhé, Suzune-chan!}''

Shion thì ngồi thẳng lưng như một vị nữ vương trên ngai vàng bóng tối của mình. Dù vẫn giữ vẻ ngoài kiêu hãnh đặc trưng, nhưng ánh mắt của cậu ấy lại thể hiện một sự nghiêm túc tuyệt đối. Mỗi khi hoàn thành một câu hỏi, tôi thấy cậu ấy lại khẽ nhếch mép cười, như thể vừa chinh phục được một đạo quân kẻ thù hắc ám nào đó.

Yae và Honoka cũng không hề kém cạnh. Ánh mắt của Yae mang theo sự quyết liệt của một vận động viên trên đường chạy marathon, còn Honoka thì cẩn thận rà soát từng dòng chữ trong bài thi Lịch sử với một sự tỉ mỉ chưa từng thấy. Tôi cảm thấy tự hào vô cùng khi thấy những nỗ lực của chúng tôi trong những ngày qua đã bắt đầu đơm hoa kết trái.

Hội ``báo động đỏ'' đang thực sự chiến đấu. Họ chiến đấu không chỉ cho điểm số của bản thân, mà còn cho lời hứa cùng nhau đi dự lễ hội, và trên hết là để chứng minh rằng nỗ lực của Naomi và nhóm gia sư là không hề uổng phí.

Trải qua những ngày thi liên tiếp, từ các môn tự nhiên hóc búa đến các môn xã hội dài dằng dặc, tinh thần của lớp 1-C vẫn không hề sụt giảm. Khi hồi chuông kết thúc môn thi cuối cùng vang lên, một sự tĩnh lặng bao trùm trong giây lát, trước khi toàn bộ học sinh đồng loạt buông bút và thở phào nhẹ nhõm. Cuộc chiến tri thức khốc liệt cuối cùng cũng đã khép lại.

Sau khi hồi chuông kết thúc môn thi cuối cùng vang lên, cả lớp chúng tôi dường như trút bỏ được một gánh nặng nghìn cân. Những tiếng hò reo, tiếng thở phào và cả những tiếng bàn tán xôn xao về đề thi tràn ngập khắp các hành lang. Tuy nhiên, niềm vui ấy vẫn chưa thực sự trọn vẹn khi kết quả chính thức vẫn còn là một ẩn số.

Để giải tỏa sự thấp thỏm, cả nhóm chúng tôi quyết định tụ tập lại sau giờ học để đối chiếu đáp án. Kuon và chị Kaigou, với sự nhạy bén và kiến thức vững chắc, được mọi người tin tưởng giao nhiệm vụ ``chấm thử'' dựa trên những gì các bạn còn nhớ và nháp lại.

Thế nhưng, ngay khi bắt đầu, tôi đã nhận thấy một sự đối lập rõ rệt trong phong cách của hai người khiến bầu không khí trở nên căng thẳng một cách thú vị.

Onii-sama ngồi ở bàn đầu, gương mặt lạnh lùng như một vị thẩm phán đang chuẩn bị tuyên án. Anh ấy cầm bút đỏ, rà soát từng dòng, từng chữ một cách tàn nhẫn. Với anh, sự chính xác là tuyệt đối. Những câu trắc nghiệm mà các bạn bảo là ``phân vân giữa A và B rồi chọn đại'' đều bị anh ấy thẳng tay đánh một dấu V\footnote{Ở Nhật Bản, khi chấm bài dấu V thường mang nghĩa là sai, trong khi chữ O nghĩa là đúng và $\Delta$ (hình tam giác) mang nghĩa là đúng một phần.} đỏ chót.

``Trong khoa học, không có khái niệm `may mắn','' anh lạnh nhạt nói khi Suzune đang mếu máo vì mất điểm một câu lý thuyết. ``Nếu cậu không chắc chắn, nghĩa là cậu chưa nắm vững. Và trong tự luận, sai kết quả cuối cùng nghĩa là toàn bộ logic của cậu có vấn đề. Sai là sai.''

Tôi đứng cạnh, nhìn những dấu V và những đường gạch đỏ rực trên bài của cô bạn mà không khỏi xót xa. Tôi hiểu rằng Onii-sama đang cố tình chấm thật chặt để chuẩn bị cho kịch bản tệ nhất, giúp các bạn không bị sốc nếu điểm chính thức có thấp. Nhưng nhìn cách anh ấy ``tàn sát'' hy vọng của mọi người, tôi vẫn không kìm được mà thì thầm: ``\hl{Onii-sama, anh có cần phải khắc nghiệt thế không? Nhìn Suzune-chan sắp khóc tới nơi rồi kìa.}''

Tôi chợt nhớ lại một lần Onii-sama từng nhắc về người em gái của mình ở thế giới của ảnh - Ryuuhou. Dù tôi không rõ cô ấy bao nhiêu tuổi, nhưng qua lời anh ấy kể, tôi biết anh cũng từng chấm bài cho cô ấy với thái độ tương tự. Anh ấy bảo rằng ngay cả với người thân nhất, anh cũng tuyệt nhiên không nương tay dù chỉ nửa điểm. Với anh, tri thức là sự nghiêm túc, và việc bao che cho cái sai của người thân chính là sự tàn nhẫn nhất đối với tương lai của họ. Nghĩ đến đó, tôi vừa thấy nể phục sự chính trực của anh, vừa thấy mừng vì mình không phải là học trò trực tiếp của ``ông anh'' máu lạnh này.

Trái ngược hoàn toàn với anh ấy, Onee-sama lại mang đến một luồng gió hy vọng. Chị ngồi ở góc bàn bên kia, vừa chấm vừa khẽ gật đầu khích lệ mọi người. Dù phần trắc nghiệm chị cũng khá nghiêm túc, nhưng ở phần tự luận, chị lại vô cùng hào phóng.

``Ý tưởng này tốt đấy Suzune-san,'' chị tôi mỉm cười, đánh một hình tam giác vào lề vở. ``Dù kết quả cuối cùng bị nhầm một dấu trừ, nhưng cách cậu thiết lập phương trình hoàn toàn đúng. Tôi vẫn tính điểm vớt cho phần tư duy này.'' Chị chấp nhận mọi cách giải, miễn là nó logic và dẫn đến cùng một bản chất vấn đề. Sự ``thoáng'' này của Onee-sama khiến nhóm báo động đỏ như được hồi sinh từ cõi chết.

Đến lượt tôi được chấm thử, đứng giữa hai người, tôi có cảm giác như đang bị kẹt giữa hai luồng tư duy trái ngược hoàn toàn. Nhìn sang Onii-sama, người vẫn đang miệt mài gạch những đường đỏ sắc lẹm, và những cái lắc đầu nhẹ, tôi vẫn có một chút gì đó hãi hùng. Anh ấy thực sự không khoan dung với những câu trả lời không đúng theo cách giải trong sách giáo khoa.

Trong khi đó, Onee-sama lại dành cho tôi một cái nhìn dịu dàng hơn hẳn. Mỗi khi tôi làm sai một ý nhỏ, thay vì gạch bỏ hoàn toàn, chị tôi dùng bút chì khoanh lại và kiên nhẫn giải thích cho em hiểu bản chất. Tôi cảm nhận được sự thấu hiểu của một người chị ruột luôn muốn em mình tìm thấy niềm vui trong việc học thay vì nỗi sợ hãi trước những con số.

Tôi cảm thấy may mắn vì có một người chị như Kaigou-onee-sama để làm dịu đi sự khắc nghiệt của ``ông anh'' Kuon-onii-sama, nhưng đồng thời cũng thầm cảm ơn sự cứng rắn của anh ấy, vì nó đã rèn cho tôi một bản lĩnh vững vàng trước mọi áp lực.

Phản ứng của mọi người trước hai ``thái cực'' chấm điểm này thật sự đa dạng. Suzune nhìn vào con số 40 của anh tôi rồi lại nhìn con số 60 của chị tôi, gương mặt cô nàng biến đổi liên tục từ tái mét sang hồng hào: ``Mình cảm giác như đang đi tàu lượn siêu tốc vậy! Một bên là địa ngục, một bên là thiên đường!''

Yae và Honoka cũng chung cảnh ngộ. Cô nàng Thể chất thì vò đầu bứt tai trước sự khắt khe của anh nhưng lại nhảy cẫng lên khi Kaigou công nhận cách giải ``cơ bắp'' của cô. Cô nàng bờm Cáo sa mạc thì liên tục chắp tay cầu nguyện mỗi khi Kaigou lật sang trang mới.

Inari và Kanade đứng cạnh, dù điểm số của họ khá an toàn nhưng cũng không khỏi hồi hộp thay cho các bạn. Koishin, cô nàng tsundere, dù miệng vẫn lẩm bẩm rằng ``chấm thế nào chẳng được'', nhưng đôi mắt vẫn không rời khỏi cây bút đỏ của hai người, đôi tay khẽ siết chặt vạt áo mỗi khi Onii-sama cau mày.

Hội bạn nòng cốt cũng bắt đầu đưa ra ý kiến. Uzuki và Azuki thì nghiêng về phía chị tôi, họ cho rằng sự động viên tinh thần lúc này là quan trọng nhất. Nanase thì đồng tình với cách chấm của anh tôi hơn, cho rằng: ``Phải cứng rắn như thế mới đúng. Thà chuẩn bị tâm lý cho điều tồi tệ nhất còn hơn là hy vọng hão huyền.''

Touka nhìn hai người chấm bài, ánh mắt vẫn giữ vẻ chuyên nghiệp: ``Cách chấm của Kuon-san giúp mọi người không chủ quan, còn Kaigou-san thì giúp mọi người không bỏ cuộc. Cả hai đều cần thiết.''

Wata khẽ mỉm cười, giọng nói dịu dàng như gió thoảng: ``Dù kết quả thế nào, sự nỗ lực của mọi người trong những ngày qua mới là điều quý giá nhất. Thần linh chắc chắn sẽ thấu hiểu mà.''

``Này, câu đó phải là tôi nói mới đúng đó,'' Sakura khẽ nhíu mày khi cô bạn bờm Cừu nhắc tới Thần linh.

Naomi ngồi im lặng trên chiếc ghế cao, đôi mắt lục bảo dõi theo từng sắc thái cảm xúc đang nhảy múa trên gương mặt các anh chị. Em nhìn thấy nỗi sợ hãi Onii-sama gạch bài, sự nhẹ nhõm khi Onee-sama mỉm cười, và cả sự gắn kết kỳ lạ giữa họ. Em không nói gì, nhưng một vầng sáng dịu nhẹ lại khẽ tỏa ra, giữ cho bầu không khí dù căng thẳng nhưng vẫn tràn đầy hơi ấm và niềm hy vọng. Em nhận ra rằng, con người thật kỳ lạ, họ có thể đau khổ vì một con số, nhưng cũng có thể hạnh phúc tột cùng chỉ vì một lời công nhận nhỏ nhoi.

\begin{center}
  \uline{Vài ngày sau.}
\end{center}

Những ngày cuối tháng Sáu trôi qua trong một bầu không khí đặc quánh sự hồi hộp. Cái nóng oi ả của mùa hè dường như càng bị đốt cháy thêm bởi nỗi lo âu đang hiện rõ trên từng gương mặt trong lớp 1-C. Mỗi khi nghe thấy tiếng bước chân đi dọc hành lang, cả lớp chúng tôi lại đồng loạt nín thở, đôi mắt đổ dồn về phía cánh cửa.

Hội ``báo động đỏ'' là những người đứng ngồi không yên nhất. Suzune cứ chốc chốc lại lẩm bẩm cầu nguyện, tay vân vê tà áo đến mức nhăn nhúm. Yae ngồi rung đùi liên tục, còn Honoka thì cứ nhìn chằm chằm vào chiếc điện thoại như thể đang đợi một thông báo định mệnh nào đó. Ngay cả Shion, người vốn luôn giữ vẻ kiêu ngạo, giờ đây cũng cầm ngược cuốn sách ma thuật, đôi mắt tím thi thoảng lại liếc trộm về phía bục giảng.

``\vie{Naomi-chan, em thấy mọi người có lạ không?}'' tôi khẽ hỏi Naomi khi thấy em cứ nghiêng đầu nhìn quanh.

Cô bé gật đầu, thì thầm đáp lại: ``\vie{Mọi người\dots\ đều đang tỏa ra những luồng năng lượng rất hỗn loạn ạ. Có cả sợ hãi, hy vọng và cả\dots\ sự tuyệt vọng nữa ạ.}''

``\vie{Đó là tâm trạng của những người đang đứng trước cửa tử đó em,}'' Onii-sama chen vào, mắt vẫn không rời khỏi cửa lớp. ``\vie{Hy vọng là công sức của anh em mình mấy ngày qua không đổ sông đổ biển.}''

Ở góc cửa sổ, Mitsuki đang lo lắng đan hai tay vào nhau: ``Hotaru-chan, cậu nghĩ đề toán vừa rồi có tính điểm quá trình không? Tớ lo cái phần đạo hàm đó quá.'' Hotaru thì chỉ khẽ mỉm cười trấn an, dù tay cô ấy cũng đang nắm chặt cây bút chì. Trong khi đó, Yuki thì thản nhiên tựa cằm nhìn ra sân trường: ``Lo cũng chẳng ích gì đâu. Con số đã được viết ra rồi, chúng ta chỉ việc đón nhận định mệnh thôi.''

Nhóm gia sư thì có vẻ vững tâm hơn. Sakura vẫn không ngừng kiểm tra lại các đáp án trên tờ nháp để sẵn sàng đối chiếu cho các bạn. Nanase thì làm động tác khởi động tay chân: ``Hy vọng mọi thứ sẽ diễn ra suôn sẻ.'' Touka chỉ đáp lại với một nụ cười: ``Mọi người đã nỗ lực hết mình. Kết quả sẽ phản ánh đúng điều đó.''

Căng thẳng nhất chính là hội báo động đỏ. Suzune lẩm bẩm: ``Bánh bao\dots\ bánh bao\dots\ xin hãy phù hộ cho con\dots\'' Shion thì đang vẽ một vòng tròn nhỏ lên bàn: ``Hỡi các thế lực của vực thẳm, hãy biến điểm 4 thành điểm 6, biến dấu V thành dấu O\dots\'' Yae và Honoka thì cứ nắm tay nhau, mắt nhắm nghiền như đang chờ đợi một phép màu.

Cánh cửa lớp bỗng dưng mở toang. Yuko-sensei bước vào với một xấp bài kiểm tra dày cộm trên tay. Gương mặt cô không biểu lộ chút cảm xúc nào, đôi mắt sắc sảo nhìn lướt qua cả lớp một lượt khiến không khí lạnh lẽo bao trùm ngay tức khắc.

``Chào cả lớp,'' Yuko-sensei đặt xấp bài lên bàn, tạo nên một tiếng *BỘP* khô khốc. ``Tôi đã có kết quả cho bài thi cuối kỳ này. Như đã thông báo từ trước, chỉ một bài thi dưới 50 sẽ đồng nghĩa với một tuần lễ hội đáng nhớ\dots\ ngay tại phòng học này.''

Tiếng nuốt nước bọt ực một cái vang lên rõ mồn một giữa sự tĩnh lặng. Bầu không khí trong lớp ngay lập tức căng ra như dây đàn.

``Đầu tiên là nhà Hikawa,'' cô giáo bắt đầu đọc điểm. ``Kuon-san đạt điểm tối đa môn Toán và Tiếng Anh. Kaigou-san cũng đạt 100 điểm Toán và được điểm tối đa phần Vật lý của Khoa học tự nhiên, Tiếng Anh cũng rất cao. Suzuki-san, được điểm tối đa phần Hóa học. Môn Ngữ văn và Tiếng Anh của em cũng nằm trong nhóm dẫn đầu.''

Tôi thở phào nhẹ nhõm, khẽ siết lấy tay Onee-sama.

``\vie{Thấy chưa Naomi-chan, sự nỗ lực luôn được đền đáp mà,}'' Onee-sama xoa đầu em, ánh mắt rạng rỡ.

Naomi ngước nhìn Onii-sama, khẽ hỏi: ``\vie{Vậy còn các bạn khác thì sao ạ? Em thấy họ vẫn còn đang rất lo lắng.}''

Onii-sama khẽ gật đầu: ``\vie{Chúng ta cũng đang chờ đây. Hy vọng là cái `tường thành' đó sẽ bị Suzune-san và những người khác đập tan.}''

Những môn còn lại của ba anh em chúng tôi đều ở mức 70-80 điểm, một kết quả đủ để chúng tôi tự hào.

Tiếp theo là nhóm thần đồng. Touka, Sakura và Nanase đúng như dự đoán, gần như quét sạch điểm tối đa ở các môn thế mạnh của mình. Shino cũng giữ vững phong độ với mức điểm đều đặn từ 70 đến 90. Phần lớn các bạn còn lại trong lớp đều đạt mức từ trung bình khá đến giỏi.

Sakura nhảy cẫng lên, ôm chầm lấy Nanase: ``Tuyệt quá! Chúng mình đều làm được rồi!'' Nanase cười lớn, vỗ ngực đầy tự hào: ``Hah! Với kịch bản hào hùng thế này thì làm sao mà thất bại được chứ! Tôi đã bảo mà, phương pháp giảng dạy của tôi hiệu quả là đằng khác!'' Touka dù vẫn giữ vẻ nghiêm nghị nhưng khóe môi đã khẽ cong lên: ``Tốt lắm. Coi như công sức ôn tập cho mọi người không làm ảnh hưởng đến kết quả cá nhân.''

Bầu không khí càng trở nên nghẹt thở khi Yuko-sensei chuyển sang nhóm cuối cùng. Nhóm thuộc những người ở vòng nguy hiểm.

``Shiromi-san, em suýt soát qua môn Khoa học tự nhiên và Toán với điểm số vừa vặn chạm ngưỡng an toàn. Omori-san cũng vừa đủ điểm qua Khoa học xã hội và Ngữ văn.'' Ngay khi cô dứt câu, hai cô nàng hét lên một tiếng nhỏ rồi ôm chầm lấy nhau, gương mặt rạng rỡ như vừa trúng số độc đắc.

Yae hét lên, nhấc bổng Honoka xoay một vòng: ``Honoka-chan! Chúng mình sống rồi! Chúng mình không phải ở lại trường rồi!'' Cô bạn tóc vàng thì cười đến chảy nước mắt, giọng khản đặc: ``Tớ\dots\ tớ không tin được là tớ lại qua được môn Sử! Nanase-chan đúng là thánh cứu khổ cứu nạn mà!''

``Shitou-san, môn Tiếng Anh của em cũng vừa đủ điểm đỗ. Có vẻ như phép thuật bóng tối của em đã có tác dụng đúng lúc nhỉ?'' Yuko-sensei nheo mắt nhìn cô nàng vừa được nhắc tên đang thở phào nhẹ nhõm.

Shion khẽ vuốt lại mái tóc, cố lấy lại vẻ điềm tĩnh của một phù thủy hắc ám: ``Hừm, thực ra ta đã biết trước kết quả này thông qua các vì sao rồi. Nhưng\dots'' Cô nàng cúi xuống nhìn Naomi, giọng nhỏ lại nhưng đầy chân thành: ``Cảm ơn nhóc tì nhé. Nguồn linh khí của ngươi\dots\ cũng có chút giúp ích cho pháp thuật của ta.''

Cuối cùng, cô giáo nhìn về phía Suzune, người đang run cầm cập như cành cây trước gió. Yuko-sensei khẽ thở dài, vẻ mặt trở nên nghiêm trọng: ``Còn Karamomo-san\dots\ môn Toán của em chỉ đạt được 5 điểm thôi.''

Cả lớp lặng đi, toàn bộ ánh mắt hướng về cô nàng vừa bị gọi tên đó. Dẫu rằng chúng tôi không ngạc nhiên vì thành tích của cậu ấy đã quá bết bát từ trước, nhưng thảm hại đến mức này chúng tôi không hề nghĩ tới.

Suzune đứng hình trong vài giây, đôi mắt xanh lá bắt đầu nhòe lệ, miệng mếu máo: ``Năm\dots\ năm điểm sao? Mình\dots\ mình thực sự phải ở lại sao\dots?''

Tôi vội chạy lại phía cô bạn đang mếu máo kia, nắm lấy đôi vai đang run rẩy của cô ấy: ``Suzune-chan, bình tĩnh lại nào! Chỉ là một môn thôi mà, vẫn còn cách khác mà!''

Uzuki cũng lo lắng cúi xuống: ``Đúng đó, Suzune-chan, đừng khóc mà, chắc hẳn phải có nhầm lẫn gì đó chứ\dots'' Miku và Kaoru cũng quây lại, mỗi người một câu an ủi khiến bầu không khí càng thêm bi tráng.

Naomi đứng lặng người, đôi mắt lục bảo hiện rõ sự xót xa. Em định tiến lại gần để chạm vào Suzune, đôi bàn tay nhỏ nhắn khẽ đưa ra như muốn truyền đi một chút hơi ấm cuối cùng để xoa dịu nỗi đau của người chị vốn luôn tươi cười này.

Thấy cô nàng sắp khóc to thành tiếng, Yuko-sensei bỗng bật cười, tông giọng của cô đổi khác hẳn: ``Nói cho chính xác là 5 trên 10 điểm cho phần tự luận hóc búa nhất. Tổng điểm của em là 50/100 cho môn Toán. Các môn còn lại, em đều đạt điểm khá. Chúc mừng em, Karamomo-san, em đã về đích ngay sát nút!''

Suzune khựng lại, nước mắt vẫn còn đọng trên má nhưng miệng đã há hốc vì ngạc nhiên. Phải mất vài giây để bộ não của cậu ấy xử lý thông tin, rồi sau đó là một tiếng hét vang dội nhất từ đầu năm học tới giờ: ``\textbf{Mình đỗ rồi!!! Lễ hội ơi, bánh bao ơi!! Ta tới đây!!!}''

Cả lớp 1-C như nổ tung trong sung sướng. Những tiếng hò reo, tiếng vỗ tay và cả những tiếng cười vang vọng khắp các hành lang. ``Tường thành'' thi cử cuối cùng đã hoàn toàn bị phá vỡ. Không một ai bị bỏ lại phía sau.

Yuko-sensei mỉm cười hài lòng, gõ thước xuống bàn: ``Tốt lắm. Coi như các em đã vượt qua thử thách. Từ ngày mai, lớp 1-C chính thức được phép bắt tay vào việc chuẩn bị gian hàng cho Lễ hội Aogami. Tôi mong chờ sự thể hiện của các em cho lễ hội của thị trấn năm nay.''

Trong niềm vui vỡ òa đó, tất cả mọi người dường như đều nhớ tới một bóng hình nhỏ nhắn - người trở thành chỗ dựa tinh thần của cả lớp. Yae lao tới chỗ Naomi, đôi mắt sáng rực sự cảm kích: ``Naomi-chan! Tất cả là nhờ em cả đấy! Không có em giữ cho tụi chị tỉnh táo, chắc tụi chị đã gục ngã từ lâu rồi!''

Chưa kịp để ai phản ứng, cô nàng Lực điền đã bế bổng Naomi lên và tung cô bé lên cao như một người hùng chiến thắng.

``\textbf{Ôi trời, Yae-san! Cẩn thận chứ!}'' tôi hét lên, tim suýt nữa thì nhảy ra khỏi lồng ngực. Onii-sama cũng vội vàng đưa tay ra đỡ phía dưới, gương mặt vốn lạnh lùng giờ đây hiện rõ vẻ hốt hoảng: ``Này! Con bé không phải là bóng thi đấu đâu!''

Naomi, lần đầu tiên trong đời bị tung hứng như vậy, đôi mắt lục bảo mở to vì ngạc nhiên, rồi sau đó, một tiếng cười trong trẻo như chuông bạc thoát ra từ đôi môi nhỏ nhắn của em. Đó là tiếng cười thực sự đầu tiên của em kể từ khi đến với thế giới này, một tiếng cười tràn đầy sức sống và niềm hạnh phúc thuần khiết. Nhìn Naomi cười rạng rỡ giữa vòng vây của những người bạn học đang tung hô em, tôi cảm thấy một luồng ấm áp lan tỏa trong lồng ngực.

Khi tiếng cười trong trẻo của Naomi dần lắng xuống, sự tò mò lại một lần nữa kéo cả lớp về phía bục giảng, nơi Yuko-sensei vừa dán một tờ danh sách tổng hợp kết quả của cả lớp.

Honoka đứng cạnh Touka, ngước nhìn dòng điểm số cao chót vót của cô bạn rồi thở dài thườn thượt: ``Trời ạ, Touka-chan đúng là quái vật mà. Nhìn hàng điểm 100 kia kìa, tớ chỉ ước mình san sẻ được một phần mười cái bộ não thiên tài đó thôi.''

Azuki cũng gật đầu tán thưởng, đôi mắt cô lấp lánh sự ngưỡng mộ: ``Không chỉ Touka-chan đâu, cả nhóm Thần đồng đều giữ phong độ thật đáng sợ. Thật mừng là họ ở lớp mình chứ không phải đối thủ ở lớp khác.''

Trong khi đó, hai chị em tôi lại đứng quan sát biểu cảm của Suzune và Shion. Cả hai cô nàng đang ôm chầm lấy nhau, mặt mũi tèm lem nước mắt nước mũi nhưng miệng thì cười ngoác tới tận mang tai.

``\vie{Em thấy đấy, Suzu,}'' Onee-sama khẽ huých tay tôi, chỉ về phía điểm 50 của hai người họ. ``\vie{Đó không còn là học tập nữa, đó là một màn trình diễn ảo thuật đỉnh cao. Chỉ cần thiếu một phần tư điểm thôi là giờ này họ đang ngồi viết bản kiểm điểm với Yuko-sensei rồi.}''

Đúng lúc ấy, Onii-sama bất ngờ cất giọng đầy ngạc nhiên: ``\vie{Suzuki, em đứng thứ 10 toàn lớp á? Dù kiến thức ở đây là lớp 10, vượt tận hai lớp cơ mà?}''

Tôi khẽ gãi má, cười trừ: ``\vie{À, em cũng chỉ học thêm chút ít thôi ạ. May mắn là đề năm nay trúng tủ.}'' Nghe vậy, anh cười nhếch mép, đôi mắt đỏ lấp lánh trêu chọc: ``\vie{Nhưng mà tới mức vượt cả Kaigou-chan thế này\dots\ có khi em còn cho chị mình hít khói đấy.}''

Onee-sama nghe vậy liền phồng má, giọng bực dọc: ``\vie{Ý cậu là sao hả, Kuon-kun!? Tôi chỉ kém Suzu có bảy điểm thôi đó!}''

Tôi cười xòa, vội xua tay, nghiêm túc đáp: ``\vie{Không có gì quá to tát đâu ạ. Em vẫn còn phải học nhiều lắm. Chỉ là lần này may mắn thôi.}''

Tiếng xôn xao, tranh luận và cả những lời trêu chọc về điểm số cứ thế kéo dài khắp gian phòng, biến bầu không khí căng thẳng của kỳ thi thành một buổi tiệc ăn mừng tự phát.

Chúng tôi đã cùng nhau vượt qua một thử thách, và phía trước chúng tôi là ánh sáng rực rỡ của ngày hội lớn nhất năm. Lớp 1-C đã sẵn sàng, và quan trọng nhất, Naomi đã thực sự trở thành một phần không thể thiếu của gia đình này.

\pov{Hikawa Kuon}

Ánh nắng vàng vọt của buổi hoàng hôn cuối tháng Sáu trải dài trên những dãy hành lang vắng lặng của trường trung học Aogami. Sau những màn ăn mừng náo nhiệt, phần lớn học sinh đã ra về để chuẩn bị cho kỳ nghỉ và lễ hội sắp tới.

Tôi đứng tựa lưng vào lan can sân thượng, đôi mắt dõi theo bóng dáng Kaigou và Suzuki đang dẫn Naomi ra cổng trường. Cô chị lớn khẽ quay đầu lại, ánh mắt thoáng chút ưu tư nhưng rồi lập tức gật đầu như một lời khích lệ ngầm. Suzuki cũng nhìn tôi, bàn tay em siết nhẹ quai cặp. Họ đều mơ hồ đoán được lý do tôi chọn ở lại lúc này. Naomi thì vẫn thản nhiên, em vẫy tay chào tôi với nụ cười còn đọng lại từ màn tung hứng ban nãy.

``\vie{Gặp lại sau nhé, Onii-chan.}''

Tiếng bước chân chậm rãi vang lên phía sau. Tôi không cần quay lại cũng biết đó là ai.

Nanase, Inari, Ayame và Wata - những người mà tôi coi là nòng cốt, là những người tỉnh táo nhất của lớp 1-C này. Họ đứng đó, lặng lẽ nhìn tôi như thể đã chuẩn bị tâm lý cho một điều gì đó nghiêm trọng sắp được tiết lộ.

``Kuon-san, cậu gọi riêng chúng mình lại\dots\ chắc không chỉ để ngắm hoàng hôn đâu nhỉ?'' Nanase lên tiếng. Đôi mắt xanh của cô xoáy sâu vào tôi, chờ đợi một lời giải thích từ cuộc gặp có phần đường đột này.

Tôi hít một hơi thật sâu, bàn tay vô thức chạm vào chiếc đồng hồ đeo tay - nơi thực thể Game vẫn đang ẩn mình tĩnh lặng.

``Mọi người đã bao giờ tự hỏi\dots\ tại sao gia đình tôi lại xuất hiện ở thị trấn này một cách đột ngột như vậy chưa?'' tôi bắt đầu, giọng nói phẳng lặng như mặt hồ không gợn sóng. ``Có lẽ tới đây, tôi phải tiết lộ với mọi người một bí mật rồi.''

``Cậu cứ ngập ngừng mãi thế, Kuon-san,'' Inari khẽ thu mình lại, ánh mắt cô sắc sảo như muốn thấu thị tâm can tôi. ``Nếu đã giữ chúng tôi lại đây vào giờ này, hẳn đó không phải là chuyện có thể nói qua loa đâu.''

Ayame cũng gật đầu, vẻ mặt nghiêm túc: ``Đúng vậy. Dù là chuyện gì, chúng tôi cũng sẵn sàng lắng nghe. Cậu đừng giữ một mình nữa.''

Tôi nhìn họ, khẽ thở dài: ``Thực sự thì tôi cũng không muốn phải nói ra đâu, nhưng\dots\ không sớm thì muộn thì điều đó phải xảy ra thôi. Tôi muốn mọi người nghe chính miệng tôi nói trước khi những tin đồn không hay bắt đầu lan rộng.''

``Vậy thì hãy nói đi,'' Hội trưởng Hội học sinh thúc giục, bước tới một bước. ``Chúng tôi nghe đây.''

Tôi thở một hơi dài, đôi mắt đỏ nhìn thẳng vào họ. ``Chuyện là\dots\ Không chỉ tôi, Kaigou-chan hay Suzuki\dots\ mà ngay cả thực thể đang trú ngụ trong chiếc đồng hồ này - thứ mà chúng tôi gọi là Game - tất cả chúng tôi \emph{đều không thuộc về thế giới này}.''

Một sự im lặng bao trùm. Inari nghiêng đầu, đôi mắt cô hiện lên sự bàng hoàng nhưng không quá ngạc nhiên. Ayame và Wata thì nhìn nhau, đôi môi khẽ mím lại.

``Chúng tôi đến từ một nơi rất xa, bị cuốn vào đây bởi những sự cố ngoài ý muốn,'' tôi tiếp tục, ánh mắt hướng về phía đường chân trời đang nhuốm màu đỏ sẫm. ``Và có một điều tôi cần mọi người biết trước khi Lễ hội Aogami bắt đầu. Khi lời nguyền ở thị trấn này được hóa giải hoàn toàn\dots\ khi mục đích của chúng tôi ở đây kết thúc\dots\ gia đình Hikawa sẽ rời đi.''

``Rời đi sao?'' Hội phó Hội học sinh thốt lên, giọng cô run rẩy. ``Ý cậu là\dots mãi mãi?''

Tôi im lặng một lúc lâu trước khi khẽ gật đầu: ``Có lẽ vậy. Chúng tôi không chắc về ngày tái ngộ, thậm chí là có cơ hội gặp lại nhau hay không. Vì vậy, tôi muốn nói rõ điều này với những người mà tôi tin tưởng nhất.''

``Kuon-san nói đúng đấy.''

Một giọng nói quen thuộc vang lên từ phía cầu thang. Shino và Sakura bước ra từ bóng tối của hành lang. Gương mặt họ bình thản, như thể đã biết rõ mọi chuyện từ lâu.

``Shino-san? Sakura-san?'' Ayame ngạc nhiên hỏi. ``Tại sao các cậu lại\dots''

Shino bước tới cạnh tôi, nụ cười của cô ấy giờ đây mang một chút đượm buồn: ``Đúng vậy. Bọn mình đã biết chuyện này kể từ khi Kuon-san và những người đi cùng cậu đặt chân vào lớp 1-C. Sự hiện diện của họ\dots\ thực chất là một phép màu ngắn ngủi dành cho Aogami.''

Sakura khoanh tay, đôi mắt cô nhìn xa xăm: ``Họ là những người khách lữ hành qua đường, nhưng lại là những người sẽ giúp chúng ta hóa giải xiềng xích của quá khứ. Một khi thị trấn này được tự do, họ cũng sẽ phải trở về nơi họ thuộc về.''

Tôi im lặng nhìn cô nàng Vu nữ, rồi nhìn sang những người còn lại. Sự thật này vốn dĩ đã được tôi cân nhắc rất kỹ trước khi nói ra, và đây có lẽ là thời điểm thích hợp nhất để tiết lộ nó.

``Đó cũng là lý do\dots'' tôi khẽ nói, giọng đượm chút xa xăm. ``Vì sao vào ngày Lễ tình nhân vừa rồi, tôi đã không chọn bất kỳ bức thư tỏ tình nào. Không phải vì tôi không trân trọng tình cảm của mọi người, mà vì ngay từ lúc đó, tôi đã biết rõ mình chỉ là một kẻ lữ hành qua đường. Việc gieo rắc một mầm mống hy vọng vào một tương lai không có tôi\dots\ đó mới thực sự là điều tàn nhẫn nhất.''

Mọi người lặng đi. Câu trả lời cho sự lạnh lùng hồi tháng Hai cuối cùng cũng đã sáng tỏ. Tôi không muốn để lại một vết sẹo trong tim bất kỳ ai sau khi mình biến mất.

``\dots Tôi hiểu rồi.''

Nanase đứng lặng đi một lúc lâu. Cô ấy không hét lên, không khóc lóc, chỉ đơn giản là hít một hơi thật sâu để nén lại sự xúc động đang dâng trào. Cô bước tới trước mặt tôi, đôi mắt nhìn thẳng vào tôi với một sự kiên định lạ lùng.

``Kuon-san, cậu đúng là một kẻ tồi tệ khi nói ra điều này ngay lúc chúng ta vừa vượt qua kỳ thi,'' cô nàng Tiểu thư nói, giọng cô khàn đi. ``Nhưng\dots\ cảm ơn cậu vì đã thẳng thắn. Tôi ghét nhất là sự mập mờ.''

Cô dừng lại một chút, rồi tiếp lời với một nụ cười nhạt: ``Nếu vậy, lễ hội lần này không chỉ là lễ hội của thị trấn nữa. Nó sẽ là lễ hội đầu tiên và cuối cùng của lớp 1-C với đầy đủ thành viên. Chúng tôi sẽ không để các cậu rời đi với những ký ức buồn bã đâu. Chúng tôi sẽ khiến cậu phải hối hận vì đã bỏ lại một thị trấn tuyệt vời như thế này.''

Inari khẽ mỉm cười, lau đi giọt nước mắt vừa chớm trên khóe mi: ``Nanase-san đã nói vậy, thì với tư cách là Lớp trưởng, mình cũng không có ý kiến gì khác. Chúng ta sẽ làm cho lễ hội năm nay trở thành một kỷ niệm đáng nhớ không chỉ với lớp chúng ta, mà cho cả thị trấn Aogami này.''

Ayame nắm chặt hai bàn tay, ánh mắt bùng cháy một quyết tâm hiếm thấy: ``Nanase-san nói đúng. Chúng ta không có thời gian để buồn bã đâu. Nếu đây là mùa hè cuối cùng chúng ta ở bên nhau, vậy thì hãy khiến nó trở thành mùa hè rực rỡ nhất trong lịch sử Aogami này!''

Wata cũng khẽ gật đầu, nụ cười dịu dàng thường ngày giờ đây mang một chút vẻ kiên định: ``Mình sẽ dốc hết sức cho gian hàng của lớp. Chúng ta sẽ cùng nhau tạo nên những ký ức mà dù có đi đến tận cùng thế giới, cậu cũng tuyệt đối không thể nào quên được.''

Nhìn những người bạn của mình, tôi cảm thấy một gánh nặng trong lòng như được trút bỏ, nhưng thay vào đó là một nỗi buồn man mác khó tả. Hoàng hôn đã tắt hẳn, nhường chỗ cho màn đêm với những ánh sao bắt đầu lấp lánh.

Đêm đang đến, báo hiệu cho một lễ hội đã được mong đợi từ lâu. Một lễ hội của niềm vui, của sự tự do, và có lẽ là cả của những lời chia ly được báo trước.

\newpage

\trva

Dưới đây là điểm tổng kết của lớp 1-C sau kỳ thi vừa rồi:

\begin{center}
  % TO-DO: thêm hai cột "Tổng điểm (viết dạng \shortstack{Tổng \\ điểm}" và "Xếp hạng (cũng shortstack)" (cả hai đều căn giữa). Đảm bảo rằng Touka đứng đầu lớp.
  \begin{tabular}{|>{\centering\arraybackslash}p{0.05\linewidth}|>{\centering\arraybackslash}p{0.3\linewidth}|c|c|c|c|c|c|c|}
    \hline
    \multirow{2}{*}{\shortstack{STT}} & \multirow{2}{*}{\shortstack{Tên đầy đủ}} & \multicolumn{5}{c|}{Điểm tổng kết} & \multirow{2}{*}{\shortstack{Tổng \\ điểm}} & \multirow{2}{*}{\shortstack{Xếp \\ hạng}}                                                    \\ \cline{3-7}
                                      &                                          & \shortstack{Toán}                  & \shortstack{Văn}                           & \shortstack{Anh}                          & \shortstack{KHTN} & \shortstack{KHXH} &     &    \\ \hline
    1                                 & Inukai Akane                             & 85                                 & 82                                         & 88                                        & 79                & 81                & 415 & 16 \\ \hline
    2                                 & Mizuno Aku                               & 75                                 & 78                                         & 80                                        & 72                & 70                & 375 & 28 \\ \hline
    3                                 & Utachi Azuki                             & 80                                 & 85                                         & 82                                        & 78                & 75                & 400 & 20 \\ \hline
    4                                 & Onidoro Ayame                            & 90                                 & 88                                         & 92                                        & 85                & 89                & 444 & 8  \\ \hline
    5                                 & Shirayuki Inari                          & 95                                 & 92                                         & 94                                        & 88                & 90                & 459 & 4  \\ \hline
    6                                 & Mizuta Uzuki                             & 78                                 & 80                                         & 75                                        & 82                & 76                & 391 & 25 \\ \hline
    7                                 & Hikawa Kaigou                            & 100                                & 75                                         & 95                                        & 85                & 78                & 433 & 11 \\ \hline
    8                                 & Ukai Kaoru                               & 82                                 & 85                                         & 80                                        & 78                & 84                & 409 & 18 \\ \hline
    9                                 & Tokiwa Kana                              & 75                                 & 80                                         & 85                                        & 70                & 82                & 392 & 24 \\ \hline
    10                                & Amano Kanade                             & 88                                 & 90                                         & 85                                        & 82                & 86                & 431 & 12 \\ \hline
    11                                & Hikawa Kuon                              & 100                                & 78                                         & 100                                       & 80                & 94                & 452 & 6  \\ \hline
    12                                & Akaba Koishin                            & 85                                 & 90                                         & 82                                        & 75                & 88                & 420 & 15 \\ \hline
    13                                & Ibara Saki                               & 78                                 & 82                                         & 75                                        & 80                & 80                & 395 & 23 \\ \hline
    14                                & Shinomiya Sakura                         & 98                                 & 90                                         & 92                                        & 90                & 90                & 460 & 3  \\ \hline
    15                                & Kisaragi Saya                            & 80                                 & 85                                         & 82                                        & 90                & 88                & 425 & 14 \\ \hline
    16                                & Shitou Shion                             & 65                                 & 68                                         & 50                                        & 62                & 70                & 315 & 29 \\ \hline
    17                                & Fukami Shiina                            & 82                                 & 78                                         & 80                                        & 85                & 75                & 400 & 21 \\ \hline
    18                                & Misora Shino                             & 85                                 & 88                                         & 90                                        & 82                & 86                & 431 & 13 \\ \hline
    19                                & Hikawa Suzuki                            & 75                                 & 95                                         & 92                                        & 90                & 88                & 440 & 9  \\ \hline
    20                                & Karamomo Suzune                          & 50                                 & 65                                         & 60                                        & 62                & 68                & 305 & 30 \\ \hline
    21                                & Yukihara Touka                           & 100                                & 98                                         & 95                                        & 100               & 100               & 493 & 1  \\ \hline
    22                                & Furukawa Nanase                          & 88                                 & 100                                        & 98                                        & 92                & 100               & 478 & 2  \\ \hline
    23                                & Natsuno Hanabi                           & 75                                 & 80                                         & 78                                        & 72                & 85                & 390 & 26 \\ \hline
    24                                & Shiratori Hotaru                         & 82                                 & 85                                         & 88                                        & 75                & 80                & 410 & 17 \\ \hline
    25                                & Omori Honoka                             & 65                                 & 50                                         & 68                                        & 62                & 50                & 295 & 31 \\ \hline
    26                                & Akagane Mari                             & 78                                 & 82                                         & 75                                        & 80                & 85                & 400 & 22 \\ \hline
    27                                & Kamishiro Miku                           & 80                                 & 85                                         & 82                                        & 78                & 80                & 405 & 19 \\ \hline
    28                                & Himekawa Mitsuki                         & 75                                 & 80                                         & 78                                        & 72                & 85                & 390 & 27 \\ \hline
    29                                & Shiromi Yae                              & 50                                 & 65                                         & 62                                        & 50                & 68                & 295 & 32 \\ \hline
    30                                & Nekomiya Yuka                            & 90                                 & 88                                         & 95                                        & 92                & 85                & 450 & 7  \\ \hline
    31                                & Kaizuka Yuki                             & 85                                 & 92                                         & 88                                        & 80                & 95                & 440 & 10 \\ \hline
    32                                & Hitsujikado Wata                         & 92                                 & 95                                         & 90                                        & 88                & 94                & 459 & 5  \\ \hline
  \end{tabular}
\end{center}

\end{document}